The source says ``paths of length four.'' Its incidence construction
instead gives configurations with \textbf{four vertices and three
edges}:

\[
 P-M-x-L,\qquad P\ne x,\quad M\ne L.
\]

Here P,x are points and M,L are lines of the original partial linear
space. The corresponding chamber is (P,L,(x,M)). Partial linearity
implies P is not on L. Conversely every chamber yields precisely this
path. Thus the correct unconditional equivalence is between chamber
transitivity and transitivity on these oriented three-edge paths
starting at a point. The author manuscript B2 already specifies this
configuration correctly. The source's conventional four-edge reading is
a stronger, false statement; no assertion about the author's intended
convention is needed.

Take the affine plane AG(2,3) and G=ASL(2,3). It is thick and connected,
and G is flag-transitive on it. Construction3 gives a firm residually
connected geometry, as proved directly by B.AffinePlane. That proof also
establishes that G acts regularly on its chambers. Its group order and
chamber count are

\[
 |G|=3^2|SL(2,3)|=216.
\]

There are 9·4·2·3=216 oriented three-edge paths beginning at a point.
Each extends to exactly two four-edge paths, by choosing a point on the
final line other than the middle point. Both choices are distinct from
the first point because the first point does not lie on the final line.
Hence there are 432 genuine four-edge paths. A group of order 216 cannot
be transitive on them. This supplies a counterexample to the
conventional four-edge conclusion under all of Theorem1's relevant
hypotheses.

More generally, if the finite partial linear space has $v$ points,
$s$ points on each line and $t$ lines through each point, the chamber
bijection gives exactly $vt(s-1)(t-1)$ chambers: choose $P$, then $M$
through $P$, then $x\in M\setminus\{P\}$, then $L\ne M$ through $x$.
Partial linearity ensures $P\notin L$. A group transitive on these
chambers has their number as an orbit size, so this integer divides
its order. This proves the corrected divisibility assertion of
Corollary 1 of [L].

The source's own divisibility formula in Corollary1, v t(s−1)(t−1),
equals the three-edge count. It supports the precise chamber-equivalent
formulation, but does not prove historical intent about terminology.
